\subsection{二次性爱与不应期的管理}

射精后的一段时间里，男性对性刺激"暂时关机"——这就是\textbf{不应期}。很多人关心：不应期能不能缩短？一夜能不能"再来一次"？本节讲清原理与可行的做法，也提醒别被"次数"绑架。

\vspace{0.5em}
\noindent\textbf{1. 不应期是什么}

\begin{itemize}
	\item 射精后\textbf{催乳素升高}，带来短暂的性欲抑制与困倦；同时阴茎对刺激的反应性下降，这段时间即不应期。
	\item \textbf{不应期长短因人、因龄而异}：年轻人可能数分钟到数小时，随年龄增长可延长至数小时甚至更久。这是\textbf{正常的生理差异}，不是"行不行"的指标。
	\item 不应期只针对"能否再次勃起与射精"，\textbf{并不等于"不能继续亲密"}——亲吻、抚摸、为对方服务等仍可继续。
\end{itemize}

\noindent\textbf{2. "二次性爱"的常见期待与现实}

\begin{itemize}
	\item \textbf{时长错觉}：色情片常呈现"立刻再来"，但那是剪辑与多次拍摄的结果，\textbf{不代表真实常态}。
	\item \textbf{第二次往往更久}：若在不应期后再次性交，许多人反映"更持久"——这也是部分人用"先自慰一次再同房"来延长射精时间的原因（见早泄相关章节）。
	\item \textbf{别把"能不能第二次"当成能力考核}：强行追求次数容易引发\textbf{表现焦虑}，反而损害勃起与性欲。
\end{itemize}

\noindent\textbf{3. 也许能帮上忙的做法}

\begin{itemize}
	\item \textbf{睡眠与疲劳}：不应期很大程度上受疲劳影响，睡够了自然"恢复"更快。
	\item \textbf{适度运动与体能}：心肺功能与整体体能好，恢复通常更快。
	\item \textbf{别在不应期硬来}：给身体一点时间，比强行刺激更有效；对伴侣而言，"换个方式继续亲密"往往比干等更好。
	\item \textbf{药物因素}：某些药物（如 SSRI 抗抑郁药、部分降压药）会延长不应期或影响射精；若困扰明显，应与医生谈，\textbf{不要自行停药}。
\end{itemize}

\noindent\textbf{4. 什么时候该就医}

\begin{itemize}
	\item 不应期\textbf{异常延长}、或射精后出现持续\textbf{极度疲劳、情绪低落、头痛、肌肉酸痛}，可能提示性高潮后疾病综合征（POIS）等问题，应就医评估。
	\item 若"第二次"完全无法勃起并为此持续焦虑、影响关系，可咨询男科或性治疗师。
\end{itemize}

\begin{tcolorbox}[colback=pink!5!white,colframe=red!50!black,title=贴心小叮咛]
不应期是身体的"充电时间"，长短天注定，强求不来。真正的亲密不靠次数撑场面——\textbf{一次专注而放松的性爱，远胜两次勉强而焦虑的"完成任务"}。若为"能不能再来"而焦虑，焦虑本身才是更该处理的问题。
\end{tcolorbox}
